%=====================================================================
%  Quantitative Reasoning - LAB 413
%  Lab 6: Confidence Intervals & Hypothesis Testing Basics (Week 7 Recap)
%  Instructor: Subin Na
%  Beamer conversion of Lab6_1015_PPT.pptx
%  Compile: pdflatex Lab6_beamer.tex   (run twice for footline counters)
%=====================================================================
\documentclass[aspectratio=169,11pt]{beamer}

%---------------------------------------------------------------------
%  Packages
%---------------------------------------------------------------------
\usepackage[utf8]{inputenc}
\usepackage[T1]{fontenc}
\usepackage{lmodern}
\usepackage{amsmath,amssymb}
\usepackage{graphicx}
\usepackage{booktabs}
\usepackage{tcolorbox}
\tcbuselibrary{skins}

%---------------------------------------------------------------------
%  QR-Lab 413 theme  (shared across all 11 labs)
%---------------------------------------------------------------------
\usetheme{default}
\usecolortheme{default}
\useinnertheme{circles}

\definecolor{qrnavy}{RGB}{20,44,90}
\definecolor{qrblue}{RGB}{38,80,150}
\definecolor{qrgray}{RGB}{90,95,105}
\definecolor{qrlight}{RGB}{234,239,247}
\definecolor{qrred}{RGB}{190,60,60}

\setbeamercolor{structure}{fg=qrnavy}
\setbeamercolor{frametitle}{fg=white,bg=qrnavy}
\setbeamercolor{title}{fg=qrnavy}
\setbeamercolor{block title}{fg=white,bg=qrblue}
\setbeamercolor{block body}{bg=qrlight}
\setbeamercolor{itemize item}{fg=qrblue}
\setbeamercolor{itemize subitem}{fg=qrblue}
\setbeamerfont{frametitle}{series=\bfseries}
\setbeamertemplate{navigation symbols}{}

\setbeamertemplate{frametitle}{%
  \nointerlineskip
  \begin{beamercolorbox}[wd=\paperwidth,ht=3.2ex,dp=1.2ex,leftskip=0.3cm]{frametitle}%
    \strut\insertframetitle\strut
  \end{beamercolorbox}%
}

% Footline: course | lab N | slide number  (edit "Lab N" per deck)
\newcommand{\labnum}{Lab 6}
\setbeamertemplate{footline}{%
  \leavevmode%
  \hbox{%
    \begin{beamercolorbox}[wd=.5\paperwidth,ht=2.4ex,dp=1ex,leftskip=.3cm]{author in head/foot}%
      \usebeamerfont{author in head/foot}\color{qrgray}\small Quantitative Reasoning -- LAB 413%
    \end{beamercolorbox}%
    \begin{beamercolorbox}[wd=.5\paperwidth,ht=2.4ex,dp=1ex,rightskip=.3cm]{title in head/foot}%
      \usebeamerfont{title in head/foot}\color{qrgray}\small\hfill \labnum \quad\insertframenumber\,/\,\inserttotalframenumber%
    \end{beamercolorbox}}%
}

% Number badge (01, 02, ...) used on concept slides:  \numbadge{01}
\newtcbox{\numbadge}{on line,boxsep=2pt,left=4pt,right=4pt,top=2pt,bottom=2pt,
  colback=qrnavy,colframe=qrnavy,fontupper=\bfseries\color{white},arc=2pt}

% Key-concepts callout:  \begin{keybox} ... \end{keybox}   or  \begin{keybox}[Scenario] ...
\newtcolorbox{keybox}[1][Key Concepts]{colback=qrlight,colframe=qrnavy,
  fonttitle=\bfseries,title={#1},coltitle=white,arc=2pt,left=4pt,right=4pt,
  top=3pt,bottom=3pt,boxrule=0.6pt}

%---------------------------------------------------------------------
%  Title data
%---------------------------------------------------------------------
\title{\bfseries Confidence Intervals \& Hypothesis Testing}
\subtitle{Week 7 Recap \textbullet\ Lab 6}
\author{Instructor: Subin Na}
\institute{Quantitative Reasoning -- LAB 413}
\date{October 15, 2025}

\begin{document}

%---------------------------------------------------------------------
%  Title slide
%---------------------------------------------------------------------
{\setbeamertemplate{footline}{}
\begin{frame}
  \vfill
  {\color{qrgray}\small Quantitative Reasoning -- LAB 413}\\[0.6em]
  {\usebeamerfont{title}\color{qrnavy}\huge\bfseries Confidence Intervals\\[0.15em]\& Hypothesis Testing\par}
  \vspace{0.4em}
  {\color{qrblue}\large Week 7 Recap \ \textbullet\ \ Lab 6}\par
  \vspace{1.6em}
  {\color{qrgray} Instructor: \textbf{Subin Na} \hfill October 15, 2025}
  \vfill
\end{frame}
}

%---------------------------------------------------------------------
%  Today's Plan
%---------------------------------------------------------------------
\begin{frame}{Today's Plan}
  \begin{columns}[T]
    \begin{column}{0.5\textwidth}
      \textbf{\color{qrnavy}1. Quick Recap}
      \begin{itemize}
        \item Overview
        \item Confidence Interval
        \item How CIs Behave
        \item Hypothesis Testing
        \item $p$-Values \& Significance
      \end{itemize}
    \end{column}
    \begin{column}{0.5\textwidth}
      \phantom{\textbf{x}}
      \vspace{1.6em}

      \textbf{\color{qrnavy}2. RStudio Practice}
      \begin{itemize}
        \item Regression CIs on the \texttt{Prestige} data
      \end{itemize}
    \end{column}
  \end{columns}
\end{frame}

%---------------------------------------------------------------------
%  01 Overview   (RECREATE concept boxes)
%---------------------------------------------------------------------
\begin{frame}{\numbadge{01}\ \ Overview}
  \begin{itemize}
    \item \textbf{Goal:} Quantify the uncertainty behind point estimates using probability.
    \item Two types of statistical inference:
  \end{itemize}
  \vspace{0.6em}
  \begin{columns}[T]
    \begin{column}{0.48\textwidth}
      \begin{tcolorbox}[colback=white,colframe=qrnavy,arc=3pt,halign title=center,
        title=\bfseries Confidence Intervals,coltitle=white,colbacktitle=qrnavy]
        \centering estimate \textbf{population parameters}
      \end{tcolorbox}
    \end{column}
    \begin{column}{0.48\textwidth}
      \begin{tcolorbox}[colback=white,colframe=qrnavy,arc=3pt,halign title=center,
        title=\bfseries Hypothesis Tests,coltitle=white,colbacktitle=qrnavy]
        \centering assess \textbf{evidence about claims} concerning the population
      \end{tcolorbox}
    \end{column}
  \end{columns}
\end{frame}

%---------------------------------------------------------------------
%  02 Confidence Intervals (1)
%---------------------------------------------------------------------
\begin{frame}{\numbadge{02}\ \ Confidence Intervals (1)}
  \begin{keybox}
    \begin{itemize}
      \item \textbf{Point estimate:} a single best guess of the parameter
            (often incorrect because of sampling variability).
      \item \textbf{Confidence interval (CI):} a range of plausible values for
            the parameter.
    \end{itemize}
  \end{keybox}
  \vspace{0.6em}
  \textbf{\color{qrnavy}Formula}
  \begin{equation*}
    \text{CI} \;=\; \text{point estimate} \;\pm\; \underbrace{\text{margin of error}}_{\text{random variation from sampling}}
  \end{equation*}
\end{frame}

%---------------------------------------------------------------------
%  02 Confidence Intervals (2)  -- components of the formula
%---------------------------------------------------------------------
\begin{frame}{\numbadge{02}\ \ Confidence Intervals (2)}
  \begin{equation*}
    \text{CI} \;=\; \bar{x} \;\pm\; z^{*}\cdot\frac{\sigma}{\sqrt{n}}
  \end{equation*}
  \vspace{0.2em}
  \begin{itemize}
    \item $\bar{x}$ : sample mean
    \item $z^{*}$ : critical value based on our confidence level
    \item $\dfrac{\sigma}{\sqrt{n}}$ : standard error --- how much the sample mean
          varies from sample to sample
  \end{itemize}
\end{frame}

%---------------------------------------------------------------------
%  02 Confidence Intervals (3)  -- interpretation
%---------------------------------------------------------------------
\begin{frame}{\numbadge{02}\ \ Confidence Intervals (3)}
  \begin{keybox}[Interpretation]
    A \textbf{95\% CI} means: if we repeated our sampling process many times,
    about \textbf{95\% of those intervals} would contain the true population
    parameter.
  \end{keybox}
  \vspace{0.8em}
  \begin{center}
  \begin{align*}
    \text{CI}_{90}   &= \bar{x} \;\pm\; 1.645\cdot \text{SE}\\[0.3em]
    \text{CI}_{95}   &= \bar{x} \;\pm\; 1.96\cdot \text{SE}\\[0.3em]
    \text{CI}_{99.7} &= \bar{x} \;\pm\; 2.576\cdot \text{SE}
  \end{align*}
  \end{center}
\end{frame}

%---------------------------------------------------------------------
%  03 How CIs Behave
%---------------------------------------------------------------------
\begin{frame}{\numbadge{03}\ \ How CIs Behave}
  \begin{keybox}
    The width of a confidence interval is driven by the \textbf{margin of error}.
  \end{keybox}
  \vspace{0.6em}
  \textbf{\color{qrnavy}The CI gets smaller when:}
  \begin{itemize}
    \item the confidence level decreases (smaller $z^{*}$),
    \item the population SD decreases (smaller $\sigma$),
    \item the sample size increases (larger $n$).
  \end{itemize}
  \vspace{0.5em}
  \begin{center}
    $\bar{x} \;\pm\; z^{*}\cdot\dfrac{\sigma}{\sqrt{n}}$
  \end{center}
\end{frame}

%---------------------------------------------------------------------
%  Worked example -- setup
%---------------------------------------------------------------------
\begin{frame}{Worked Example: Sleep Hours}
  \begin{keybox}[Scenario]
    The director of public health in Philadelphia wants to estimate the average
    number of hours adults sleep per night. A survey of \textbf{$n = 250$}
    randomly selected residents gives a sample mean of \textbf{6.8 hours} with a
    standard deviation of \textbf{1.4 hours}.
  \end{keybox}
  \vspace{0.6em}
  \begin{center}
    Construct a \textbf{95\% confidence interval} for the true mean number of
    sleep hours among all adults in Philadelphia. What interval would you report?
  \end{center}
\end{frame}

%---------------------------------------------------------------------
%  Worked example -- solution
%---------------------------------------------------------------------
\begin{frame}{Worked Example: Solution}
  \textbf{\color{qrnavy}Given}
  \begin{equation*}
    \bar{x} = 6.8, \qquad s = 1.4, \qquad n = 250
  \end{equation*}
  \vspace{0.2em}
  \textbf{\color{qrnavy}Compute the 95\% CI}
  \begin{align*}
    \bar{x} \;\pm\; z^{*}\cdot\frac{\sigma}{\sqrt{n}}
      &= 6.8 \;\pm\; 1.96\cdot\frac{1.4}{\sqrt{250}}\\[0.3em]
      &= 6.8 \;\pm\; 1.96\cdot 0.0886
       \;=\; 6.8 \;\pm\; 0.174
  \end{align*}
  \vspace{0.1em}
  \begin{center}
  \begin{tcolorbox}[colback=qrlight,colframe=qrnavy,arc=2pt,boxrule=0.6pt,width=0.6\textwidth,halign=center]
    $\Rightarrow\; (6.63,\ 6.97)$ hours
  \end{tcolorbox}
  \end{center}
\end{frame}

%---------------------------------------------------------------------
%  04 Hypothesis Testing
%---------------------------------------------------------------------
\begin{frame}{\numbadge{04}\ \ Hypothesis Testing}
  \textbf{\color{qrnavy}The Setup}
  \begin{columns}[T]
    \begin{column}{0.5\textwidth}
      \begin{block}{Null hypothesis $H_0$}
        Status quo, ``no effect'' --- what we try to disprove.
      \end{block}
    \end{column}
    \begin{column}{0.5\textwidth}
      \begin{block}{Alternative hypothesis $H_1$}
        What we hope to prove.
      \end{block}
    \end{column}
  \end{columns}
  \vspace{0.6em}
  \textbf{\color{qrnavy}Test statistic}
  \begin{equation*}
    z \;=\; \frac{\text{estimate} - \text{hypothesized value}}{\text{standard deviation of the estimate}}
      \;=\; \frac{\text{observed} - \text{null guess}}{\text{SE}}
  \end{equation*}
  \begin{center}
    Large $|z| \;\Rightarrow\;$ evidence against $H_0$.
  \end{center}
\end{frame}

%---------------------------------------------------------------------
%  05 P-Values & Significance
%---------------------------------------------------------------------
\begin{frame}{\numbadge{05}\ \ $p$-Values \& Significance}
  \begin{keybox}[$p$-value]
    The probability of observing data as extreme as ours \emph{if $H_0$ were true}.
  \end{keybox}
  \vspace{0.4em}
  \begin{columns}[T]
    \begin{column}{0.54\textwidth}
      \small
      \renewcommand{\arraystretch}{1.25}
      \begin{tabular}{@{}lll@{}}
        \toprule
        \textbf{$p$-value range} & \textbf{Interpretation} & \textbf{Symbol}\\
        \midrule
        $p \geq 0.10$          & Not significant        &      \\
        $0.05 \leq p < 0.10$   & Marginally significant & $*$  \\
        $p < 0.05$             & Significant            & $**$ \\
        $p < 0.01$             & Highly significant     & $*\!*\!*$\\
        \bottomrule
      \end{tabular}
    \end{column}
    \begin{column}{0.44\textwidth}
      \textbf{\color{qrnavy}Decision rule ($\alpha$)}
      \begin{itemize}
        \item If $p < \alpha$: reject $H_0$.
        \item If $p \geq \alpha$: fail to reject $H_0$.
      \end{itemize}
      \vspace{0.4em}
      \begin{tcolorbox}[colback=qrlight,colframe=qrred,arc=2pt,boxrule=0.6pt]
        Never ``accept $H_0$'' --- only \textbf{reject} or \textbf{fail to reject}.
      \end{tcolorbox}
    \end{column}
  \end{columns}
\end{frame}

%---------------------------------------------------------------------
%  Transition to RStudio
%---------------------------------------------------------------------
{\setbeamertemplate{footline}{}
\begin{frame}
  \vfill
  \centering
  {\color{qrnavy}\Huge\bfseries RStudio Practice}\\[0.8em]
  {\color{qrblue}\large Regression Confidence Intervals with the \texttt{Prestige} data}\\[0.4em]
  {\color{qrgray} See \texttt{Lab6.Rmd} / \texttt{Lab6\_1015.R}}
  \vfill
\end{frame}
}

\end{document}
